% Supporting Information
% Scientific definitions, generated closure-table values, and numerical results remain unchanged.

\documentclass[11pt,a4paper]{article}

\usepackage[margin=2.0cm]{geometry}
\usepackage[T1]{fontenc}
\usepackage[utf8]{inputenc}
\usepackage{lmodern}
\usepackage{microtype}
\usepackage{amsmath,amssymb}
\usepackage{siunitx}
\DeclareSIUnit{\calorie}{cal}
\usepackage{booktabs,tabularx,longtable}
\usepackage{array}
\usepackage{graphicx}
\usepackage{float}
\usepackage{pdflscape}
\usepackage{caption}
\usepackage{xcolor}
\usepackage[numbers,sort&compress]{natbib}
\usepackage[colorlinks=true,linkcolor=blue!45!black,
            citecolor=green!35!black,urlcolor=blue!55!black]{hyperref}
\usepackage[nameinlink,noabbrev]{cleveref}
\crefname{figure}{Fig.}{Figs.}
\Crefname{figure}{Figure}{Figures}
\crefname{table}{Table}{Tables}
\Crefname{table}{Table}{Tables}
\crefname{section}{Section}{Sections}
\Crefname{section}{Section}{Sections}

\graphicspath{{figures/}{figures/final/}}
\setlength{\emergencystretch}{2em}
\newcommand{\papertitle}{Matched MOF Families Reveal Regime-Dependent Pore-Geometry Effects on Adsorption}

\newcommand{\sifigure}[5]{%
  \begin{figure}[H]
    \centering
    \IfFileExists{figures/#1}{%
      \includegraphics[width=0.96\linewidth]{figures/#1}%
    }{%
      \IfFileExists{figures/final/#1}{%
        \includegraphics[width=0.96\linewidth]{figures/final/#1}%
      }{%
        \fbox{\parbox[c][6.0cm][c]{0.92\linewidth}{%
          \centering\textbf{#3}\\[0.7em]
          \small #4\\[0.7em]
          \textit{Expected file:}\\[0.2em]
          \nolinkurl{figures/#1}\\[-0.1em]
          \textit{or}\\[-0.1em]
          \nolinkurl{figures/final/#1}}}%
      }%
    }
    \caption{\textbf{#3.} #5}
    \label{#2}
  \end{figure}%
}

% Supporting Information numbering
\renewcommand{\thetable}{S\arabic{table}}
\renewcommand{\thefigure}{S\arabic{figure}}

\begin{document}

\begin{center}
{\LARGE \textbf{Supporting Information}}\\[0.5em]
{\Large \textbf{\papertitle}}\\[1.2em]

{\large \textbf{Hosein Alimardani}$^{1,\dagger}$, \textbf{Shayan Abaei}$^{2,\dagger}$,
\textbf{Mehrdad Asgari}$^{3,*}$}\\[0.9em]

{\normalsize $^{1}$Faculty of Engineering, University of Tehran, Tehran, Iran}\\[0.35em]
{\normalsize $^{2}$School of Chemical Engineering, Faculty of Engineering, University of Tehran, Tehran, Iran}\\[0.35em]
{\normalsize $^{3}$Lucy Cavendish College, University of Cambridge, Cambridge CB3 0BU, UK}\\[0.8em]

{\normalsize $^{\dagger}$These authors contributed equally.}\\[0.35em]
{\normalsize $^{*}$Corresponding author:
\href{mailto:ma2000@cam.ac.uk}{\texttt{ma2000@cam.ac.uk}}}
\end{center}

\vspace{1em}
\tableofcontents
\clearpage

\section{Purpose and analysis hierarchy}
\label{sec:purpose}

This Supporting Information (SI) provides the complete evidence trail for the
matched-family analysis while keeping the inferential hierarchy explicit. The
outcome-blind construction and frozen primary comparison are documented in
\cref{sec:cohort,sec:geometry-support,sec:frozen-pairs,sec:primary};
secondary adsorption, process, concordance, pressure-transition, and archetype
results follow in
\cref{sec:secondary-adsorption,sec:secondary-process,sec:concordance,sec:transitions,sec:archetypes}. Dominant-context and database-provenance robustness
are reported in \cref{sec:robustness,sec:provenance}. The generated complete
secondary-result closure and source/code map are retained as Sections 14--15 and
\cref{tab:S23,tab:S24,tab:S25}. Post-freeze attribution, uptake-basis/density, and
same-source checks are reported in
\cref{sec:attribution,sec:basis-density,sec:source-heterogeneity}; the structural
case ledger is \cref{sec:structures}. The document closes with database-specific
limitations, the audit trail, and the interpretive boundary in
\cref{sec:limitations,sec:audit,sec:boundary}.

ARC--MOF supplies the charged CIFs, chemistry-aware representations, recorded
topology and dimensionality, pore descriptors, adsorption records, and process
metrics used here \cite{burner2023arcmof}. REPEAT is the periodic charge
methodology associated with the charged structures \cite{campana2009repeat};
high-throughput pore geometry follows established porous-material analysis
methodology \cite{willems2012zeopp}. CIF parsing used pymatgen
\cite{ong2013pymatgen}, RAC descriptors originate from the revised autocorrelation
framework \cite{janet2017rac}, and recorded topology was treated as a coarse
architecture label consistent with automated topology identification
\cite{zoubritzky2022crystalnets}.

The inferential sequence is fixed throughout this document. Recorded chemistry and
broad architecture were fixed before geometry thresholding; the 10\% geometry
criterion was frozen before deterministic pair construction; pair identity and
orientation were frozen before outcome values were read; and CH$_4$ uptake at
298 K and 65 bar in \si{\milli\mole\per\gram} was frozen as the primary endpoint
before its values were opened. All other gases, pressures, uptake bases, energetic
quantities, process metrics, concordance summaries, archetypes, source strata, and
structural case studies are secondary or post-freeze analyses. No minimum
practically meaningful primary effect was prespecified. The family bootstrap
resamples matched families and does not propagate per-framework Monte Carlo
sampling uncertainty. Secondary analyses remain descriptive; none is promoted to
primary status after value inspection.


\section{Outcome-blind cohort construction}
\label{sec:cohort}

\subsection{Stage-level census}

The outcome-blind population flow is summarized in \cref{tab:S1}. This census is the numerical basis for main-text Figure 1A; chemistry/architecture definitions are kept separate from geometry support so that the order of information access remains explicit.


\begin{table}[H]
\centering
\caption{Cohort flow through outcome-blind construction and pair freezing.}
\label{tab:S1}
\begin{tabularx}{\textwidth}{Xrr}
\toprule
Stage / population & Families or groups & Frameworks \\
\midrule
ARC--MOF registry & -- & 279,019 \\
Charged CIFs parsed successfully & -- & 242,363 \\
Complete chemistry signatures & 126,062 & 154,150 \\
Multi-member same-recorded-chemistry families & 20,278 & 48,366 \\
Validated same-recorded-architecture groups & 19,068 & 44,511 \\
Groups with measurable geometry variation & 18,390 & 41,796 \\
10\% working pore-geometry threshold & 11,072 & 26,146 \\
Frozen primary geometry pairs & 11,072 & 22,144 unique \\
\bottomrule
\end{tabularx}
\end{table}

\subsection{Chemistry and architecture definitions}

Architecture-status outcomes are reported in \cref{tab:S2}. The architecture filter is deliberately conservative: missing labels are treated as ambiguous rather than reconstructed, and geometry remains excluded from the architecture key.


The Stage 1.3 chemistry key used five components: scale-invariant whole-CIF
elemental composition, scale-invariant metal-element composition, and ARC--MOF
metal-, linker-, and functional-chemistry cluster assignments. Geometry, topology,
dimensionality, geometric-cluster labels, partial charges, and all adsorption or
process outcomes were excluded from the chemistry key. Frameworks were eligible
only after successful CIF parsing and raw-CIF element auditing, one ordered parsed
structure, at least one metal element, complete RAC evidence, and valid nonnegative
chemical-cluster labels.

Within each multi-member chemistry family, Stage 1.4 required exact agreement in
recorded topology and dimensionality. Missing architecture labels were not imputed.
No CIF reparsing, inferred bonding, CrystalNN reconstruction, or local-environment
inference was used.

\begin{table}[H]
\centering
\caption{Stage 1.4 architecture statuses.}
\label{tab:S2}
\begin{tabularx}{0.82\textwidth}{Xr}
\toprule
Status & Frameworks \\
\midrule
Validated same architecture & 44,511 \\
Ambiguous: missing topology & 2,429 \\
Architecture singleton & 1,406 \\
Ambiguous: missing dimensionality & 19 \\
Ambiguous: missing topology and dimensionality & 1 \\
\bottomrule
\end{tabularx}
\end{table}

Only 27 source chemistry families split into more than one validated architecture.

\section{Within-family geometry variation and threshold support}
\label{sec:geometry-support}

Descriptor-level variation and threshold support are reported in \cref{tab:S3,tab:S4,tab:S5}. These tables provide the numerical support behind main-text Figure 1B--C and document that the 10\% threshold was selected from outcome-blind structural support rather than adsorption performance.


All 44,511 frameworks retained after Stage 1.4 had complete values for the ten
geometry descriptors carried forward. Geometry variation was first measured
descriptively, with no scientific contrast threshold.

\begin{table}[H]
\centering
\caption{Descriptor-level within-family variation. Fractional span is
$(x_{\max}-x_{\min})/\max(|x_{\min}|,|x_{\max}|)$.}
\label{tab:S3}
\begin{tabular}{lrrrr}
\toprule
Descriptor & Families varying & Fraction & Median span & 90th pct. span \\
\midrule
UC volume & 15,570 & 0.817 & 0.0037 & 0.1301 \\
Density & 17,298 & 0.907 & 0.0050 & 0.5000 \\
ASA & 18,262 & 0.958 & 0.0571 & 0.4393 \\
AVA & 18,288 & 0.959 & 0.0721 & 0.7707 \\
AVAf & 18,264 & 0.958 & 0.0644 & 0.7458 \\
POAVA & 18,231 & 0.956 & 0.0348 & 0.4834 \\
POAVAf & 18,003 & 0.944 & 0.0301 & 0.4220 \\
$D_i$ & 18,317 & 0.961 & 0.0537 & 0.3583 \\
$D_f$ & 18,247 & 0.957 & 0.0648 & 0.4137 \\
$D_{\mathrm{if}}$ & 18,338 & 0.962 & 0.0568 & 0.3606 \\
\bottomrule
\end{tabular}
\end{table}

A total of 18,390 of 19,068 validated groups showed measurable variation in at
least one descriptor; 678 groups did not. Among the varying groups, the median
number of varying descriptors was 10 and the median strongest within-family
fractional span was 0.1353.

\begin{table}[H]
\centering
\caption{Outcome-blind support across candidate pore-geometry thresholds.}
\label{tab:S4}
\begin{tabular}{rrrr}
\toprule
Threshold & Eligible families & Fraction of all groups & Eligible frameworks \\
\midrule
5\% & 14,739 & 0.773 & 33,949 \\
10\% & 11,072 & 0.581 & 26,146 \\
20\% & 6,948 & 0.364 & 17,126 \\
30\% & 5,115 & 0.268 & 12,793 \\
\bottomrule
\end{tabular}
\end{table}

\begin{longtable}{lrrrr}
\caption{Number of architecture groups meeting each candidate threshold by
descriptor.}\label{tab:S5}\\
\toprule
Descriptor & 5\% & 10\% & 20\% & 30\% \\
\midrule
\endfirsthead
\toprule
Descriptor & 5\% & 10\% & 20\% & 30\% \\
\midrule
\endhead
ASA & 9,750 & 6,511 & 3,984 & 2,993 \\
AVA & 10,509 & 7,968 & 5,704 & 4,452 \\
AVAf & 10,112 & 7,601 & 5,294 & 4,094 \\
POAVA & 7,707 & 5,567 & 4,086 & 3,159 \\
POAVAf & 7,026 & 4,802 & 3,189 & 2,425 \\
$D_i$ & 9,483 & 5,955 & 2,927 & 2,019 \\
$D_f$ & 10,088 & 7,279 & 4,162 & 2,688 \\
$D_{\mathrm{if}}$ & 9,779 & 6,207 & 3,022 & 2,055 \\
\bottomrule
\end{longtable}

\section{Frozen geometry-pair cohort}
\label{sec:frozen-pairs}

The frozen cohort size and largest-contrast class composition are summarized in \cref{tab:S6,tab:S7}. Together with main-text Figure 1D--F, these tables show that the frozen pairs span multiple geometry modes and recorded topologies while remaining intrinsically multidimensional.


\begin{table}[H]
\centering
\caption{Summary of the frozen primary geometry-pair cohort.}
\label{tab:S6}
\begin{tabular}{lr}
\toprule
Metric & Value \\
\midrule
10\%-eligible families & 11,072 \\
Primary geometry pairs & 11,072 \\
Unique frameworks in primary pairs & 22,144 \\
Pairs also meeting 20\% threshold & 6,948 \\
Pairs also meeting 30\% threshold & 5,115 \\
Median family size & 2 \\
Families with more than two members & 2,628 \\
Median maximum pair difference & 0.2723 \\
Median mean pair difference & 0.1279 \\
\bottomrule
\end{tabular}
\end{table}

\begin{table}[H]
\centering
\caption{Pore descriptor producing the maximum relative difference in each
selected primary pair.}
\label{tab:S7}
\begin{tabular}{lrr}
\toprule
Largest-contrast descriptor & Selected pairs & Median largest contrast \\
\midrule
AVA & 3,933 & 0.3953 \\
$D_f$ & 2,211 & 0.1810 \\
ASA & 2,036 & 0.4590 \\
AVAf & 1,580 & 0.3934 \\
$D_i$ & 556 & 0.1391 \\
$D_{\mathrm{if}}$ & 540 & 0.1460 \\
POAVA & 156 & 0.1974 \\
POAVAf & 60 & 0.1855 \\
\bottomrule
\end{tabular}
\end{table}

\section{Confirmatory endpoint and primary adsorption result}
\label{sec:primary}

The primary estimate, prespecified contrast tiers, and class-resolved response are reported in \cref{tab:S8,tab:S9}. These are the complete numerical anchors for main-text Figure 2: \cref{tab:S8} reports the frozen population and stronger prespecified tiers, whereas \cref{tab:S9} separates accessible-space and diameter-labelled response classes.


Before outcome values were opened, CH$_4$ uptake at 298 K and 65 bar in
\si{\milli\mole\per\gram} was frozen as the confirmatory endpoint. The pair
orientation was frozen from lower to higher value of the already-selected
largest-contrast descriptor. The primary effect was
$\Delta q=q_{\mathrm{geometry\mbox{-}high}}-
q_{\mathrm{geometry\mbox{-}low}}$ and the population estimator was the median
across frozen families. A 10,000-replicate family bootstrap with seed 17 supplied
the percentile 95\% interval. No imputation was used.

\begin{table}[H]
\centering
\caption{Confirmatory high-pressure CH$_4$ adsorption result and prespecified
geometry-contrast sensitivity tiers.}
\label{tab:S8}
\begin{tabular}{lrrrr}
\toprule
Tier & Pairs & Median $\Delta q$ & Mean $\Delta q$ & Positive fraction \\
 & & (\si{\milli\mole\per\gram}) & (\si{\milli\mole\per\gram}) & \\
\midrule
$\geq10\%$ & 11,072 & 0.5836 & 2.4275 & 0.7050 \\
$\geq20\%$ & 6,948 & 1.6980 & 3.7880 & 0.7827 \\
$\geq30\%$ & 5,115 & 3.1645 & 5.0445 & 0.8522 \\
\bottomrule
\end{tabular}
\end{table}

For the 10\% confirmatory population, the bootstrap 95\% interval for the median
was \SIrange{0.5546}{0.6149}{\milli\mole\per\gram}. The median symmetric relative
effect was 0.0490.

\begin{table}[H]
\centering
\caption{Primary CH$_4$ response by frozen largest-contrast geometry descriptor.}
\label{tab:S9}
\begin{tabular}{lrrr}
\toprule
Largest-contrast class & Pairs & Median $\Delta q$ (\si{\milli\mole\per\gram}) &
Positive fraction \\
\midrule
ASA & 2,036 & 2.9964 & 0.9288 \\
AVA & 3,933 & 1.2201 & 0.8419 \\
POAVA & 156 & 0.5948 & 0.8077 \\
AVAf & 1,580 & 0.4646 & 0.6741 \\
POAVAf & 60 & 0.1123 & 0.5833 \\
$D_{\mathrm{if}}$ & 540 & -0.1282 & 0.4259 \\
$D_i$ & 556 & -0.1348 & 0.4263 \\
$D_f$ & 2,211 & -0.1457 & 0.4120 \\
\bottomrule
\end{tabular}
\end{table}

\section{Secondary adsorption atlas}
\label{sec:secondary-adsorption}

The complete audited \si{\milli\mole\per\gram} adsorption endpoints are listed in \cref{tab:S10}; the secondary N$_2$/H$_2$ endpoint comparison is visualized in \cref{figS:secondary-pressure}. \Cref{tab:S10} supplies the endpoint values plotted for CH$_4$/CO$_2$ in main-text Figure 3A--B and keeps the lower- and higher-pressure conditions explicit.


\begin{longtable}{llrrrr}
\caption{Secondary \si{\milli\mole\per\gram} adsorption responses using the frozen
pair orientation. The 65-bar methane-storage row is the confirmatory endpoint and
is included for reference.}\label{tab:S10}\\
\toprule
Context & Guest & Pressure & Complete pairs & Median $\Delta q$ &
Positive fraction \\
 & & (bar) & & (\si{\milli\mole\per\gram}) & \\
\midrule
\endfirsthead
\toprule
Context & Guest & Pressure & Complete pairs & Median $\Delta q$ &
Positive fraction \\
\midrule
\endhead
Methane storage & CH$_4$ & 5.8 & 11,072 & -0.0047 & 0.4947 \\
Methane storage & CH$_4$ & 65 & 11,072 & 0.5836 & 0.7050 \\
Landfill gas & CH$_4$ & 0.01 & 11,072 & -0.0004 & 0.4417 \\
Landfill gas & CH$_4$ & 4.4 & 11,072 & 0.0628 & 0.6096 \\
Landfill gas & CO$_2$ & 0.26 & 11,072 & -0.0616 & 0.3306 \\
Landfill gas & CO$_2$ & 3.2 & 11,072 & -0.0961 & 0.4378 \\
Natural-gas purification & CH$_4$ & 0.9 & 11,072 & -0.0288 & 0.4226 \\
Natural-gas purification & CH$_4$ & 9 & 11,072 & 0.1933 & 0.6143 \\
Natural-gas purification & CO$_2$ & 0.1 & 11,072 & -0.0735 & 0.3268 \\
Natural-gas purification & CO$_2$ & 1 & 11,072 & -0.1197 & 0.4184 \\
Post-combustion & CO$_2$ & 0.015 & 11,072 & -0.0179 & 0.3293 \\
Post-combustion & CO$_2$ & 0.15 & 11,072 & -0.1196 & 0.3235 \\
Post-combustion & N$_2$ & 0.075 & 11,072 & 0.0004 & 0.5424 \\
Post-combustion & N$_2$ & 0.75 & 11,072 & 0.0043 & 0.6075 \\
Pre-combustion & CO$_2$ & 0.4 & 11,072 & -0.1561 & 0.3320 \\
Pre-combustion & CO$_2$ & 16 & 11,071 & 0.5486 & 0.6690 \\
Pre-combustion & H$_2$ & 0.6 & 11,072 & 0.0026 & 0.7338 \\
Pre-combustion & H$_2$ & 24 & 11,071 & 0.0218 & 0.7054 \\
\bottomrule
\end{longtable}

\sifigure{Figure_S1.pdf}{figS:secondary-pressure}
  {Secondary N$_2$ and H$_2$ pressure responses}
  {Discrete lower- and higher-pressure median relative uptake responses for
   post-combustion N$_2$ and pre-combustion H$_2$.}
  {N$_2$ and H$_2$ were retained as secondary outcomes. Open and filled markers
   denote the available lower- and higher-pressure endpoints, respectively; the
   connecting segment is a visual link between those endpoints rather than a fitted
   pressure trajectory. Both gases show modest positive median geometry-associated
   responses, substantially smaller in absolute uptake than high-pressure CH$_4$ or
   CO$_2$.}

The full endpoint table shows that pressure response is not guest-independent.
CH$_4$ becomes more favorable at the higher stored pressure in all three methane
contexts, whereas several CO$_2$ contexts remain negative and pre-combustion
CO$_2$ changes sign. \Cref{figS:secondary-pressure} adds N$_2$ and H$_2$ without
changing that interpretation: their median responses are positive but comparatively
small. These are discrete database state points, not interpolated isotherms.

\section{Secondary process outcomes}
\label{sec:secondary-process}

Selected process metrics are summarized in \cref{tab:S11}, with the AFM--SSP trade-off shown in \cref{figS:process-tradeoff}. \Cref{tab:S11} is the numerical basis for main-text Figure 3C--D and distinguishes capacity-type objectives from separation selectivity rather than collapsing them into one performance score.


\begingroup
\small
\setlength{\tabcolsep}{3pt}
\begin{longtable}{p{0.17\textwidth}p{0.30\textwidth}rrr}
\caption{Selected process-level secondary outcomes.}\label{tab:S11}\\
\toprule
Process & Metric & Complete pairs & Median change & Positive fraction \\
\midrule
\endfirsthead
\toprule
Process & Metric & Complete pairs & Median change & Positive fraction \\
\midrule
\endhead
Landfill gas & Working capacity (\si{\milli\mole\per\gram}) & 11,008 & -0.0249 & 0.4836 \\
Landfill gas & Selectivity & 11,008 & -0.3630 & 0.3147 \\
Landfill gas & Purity & 11,008 & -0.0166 & 0.3527 \\
Landfill gas & SSP & 11,008 & -1.3663 & 0.3328 \\
Landfill gas & AFM & 11,008 & -0.6251 & 0.4947 \\
Methane storage & Working capacity (\si{\milli\mole\per\gram}) & 11,008 & 0.4668 & 0.7231 \\
Natural-gas purification & Working capacity (\si{\milli\mole\per\gram}) & 11,008 & -0.0202 & 0.4849 \\
Natural-gas purification & Selectivity & 11,008 & -0.6377 & 0.3256 \\
Natural-gas purification & Purity & 11,008 & -0.0273 & 0.3420 \\
Natural-gas purification & SSP & 11,008 & -0.7028 & 0.3352 \\
Natural-gas purification & AFM & 11,008 & -1.4204 & 0.4366 \\
Post-combustion & Working capacity (\si{\milli\mole\per\gram}) & 11,008 & -0.1028 & 0.3308 \\
Post-combustion & Selectivity & 11,008 & -4.4559 & 0.3106 \\
Post-combustion & Purity & 11,008 & -0.0153 & 0.3185 \\
Post-combustion & SSP & 11,008 & -35.1024 & 0.3150 \\
Post-combustion & AFM & 11,008 & -7.8454 & 0.3705 \\
Pre-combustion & Working capacity (\si{\milli\mole\per\gram}) & 11,008 & 0.6603 & 0.6921 \\
Pre-combustion & Selectivity & 11,008 & -4.7048 & 0.3786 \\
Pre-combustion & Purity & 11,008 & -0.0004 & 0.4584 \\
Pre-combustion & SSP & 11,008 & -155.3865 & 0.4276 \\
Pre-combustion & AFM & 11,008 & 417.4724 & 0.6578 \\
\bottomrule
\end{longtable}
\endgroup

\sifigure{Figure_S2.pdf}{figS:process-tradeoff}
  {AFM--SSP trade-offs across separation processes}
  {Median symmetric relative AFM effect versus median symmetric relative SSP effect
   for the four processes with both metrics.}
  {Process-level metrics need not respond in the same direction. The AFM--SSP
   comparison summarizes this trade-off on a common symmetric relative scale.}

The disagreement among process metrics is substantive rather than graphical.
Methane-storage and pre-combustion working capacity increase on median, while
selectivity decreases in every separation context with available values. AFM and
SSP can also occupy different response quadrants. The SI therefore reports these
metrics separately and treats them as screening objectives, not interchangeable
measures of one latent ``separation performance'' variable.

\section{Cross-condition concordance}
\label{sec:concordance}

Primary-endpoint concordance is tabulated in \cref{tab:S12}; the full endpoint-by-endpoint rank-correlation and sign-agreement matrices are shown in \cref{figS:concordance}. The fixed condition-based ordering is retained so that the matrices can be read by guest and loading regime rather than outcome-clustered after inspection.


\begin{longtable}{llrrr}
\caption{Concordance of each secondary \si{\milli\mole\per\gram} adsorption
endpoint with the frozen 65-bar CH$_4$ primary response.}\label{tab:S12}\\
\toprule
Context & Guest & Pressure (bar) & Spearman $\rho$ & Same-sign fraction \\
\midrule
\endfirsthead
\toprule
Context & Guest & Pressure (bar) & Spearman $\rho$ & Same-sign fraction \\
\midrule
\endhead
Pre-combustion & H$_2$ & 0.6 & 0.816 & 0.775 \\
Pre-combustion & CO$_2$ & 16 & 0.777 & 0.797 \\
Post-combustion & N$_2$ & 0.75 & 0.746 & 0.738 \\
Natural-gas purification & CH$_4$ & 9 & 0.640 & 0.772 \\
Landfill gas & CH$_4$ & 4.4 & 0.585 & 0.738 \\
Post-combustion & N$_2$ & 0.075 & 0.530 & 0.655 \\
Pre-combustion & H$_2$ & 24 & 0.489 & 0.608 \\
Methane storage & CH$_4$ & 5.8 & 0.289 & 0.674 \\
Landfill gas & CO$_2$ & 3.2 & 0.071 & 0.580 \\
Landfill gas & CH$_4$ & 0.01 & 0.007 & 0.541 \\
Natural-gas purification & CO$_2$ & 1 & -0.024 & 0.538 \\
Natural-gas purification & CH$_4$ & 0.9 & -0.029 & 0.545 \\
Pre-combustion & CO$_2$ & 0.4 & -0.278 & 0.448 \\
Landfill gas & CO$_2$ & 0.26 & -0.289 & 0.426 \\
Natural-gas purification & CO$_2$ & 0.1 & -0.300 & 0.428 \\
Post-combustion & CO$_2$ & 0.015 & -0.300 & 0.410 \\
Post-combustion & CO$_2$ & 0.15 & -0.318 & 0.424 \\
\bottomrule
\end{longtable}

\sifigure{Figure_S3.pdf}{figS:concordance}
  {Cross-condition effect concordance}
  {Full 18-by-18 Spearman correlation matrix and 18-by-18 sign-agreement matrix
   for the frozen pair-level \si{\milli\mole\per\gram} effects.}
  {The pair-level adsorption response is condition-dependent. Higher-loading
   conditions tend to agree more strongly with the primary 65-bar CH$_4$ response,
   whereas several low-pressure CO$_2$ conditions show weak or negative
   concordance.}

\Cref{tab:S12,figS:concordance} show that rank concordance and sign concordance are
related but not identical. Several higher-loading endpoints align strongly with the
primary methane response, whereas multiple low-pressure CO$_2$ endpoints have
negative rank correlation and below-half sign agreement. This is the pair-level
evidence underlying the operating-regime interpretation in the main text.

\section{Low-to-high pressure sign transitions}
\label{sec:transitions}

Discrete lower-to-higher-pressure sign transitions are reported in \cref{tab:S13}.


\begingroup
\small
\setlength{\tabcolsep}{3.5pt}
\begin{longtable}{p{0.20\textwidth}lrrrrr}
\caption{Pair-level sign transitions from the lower to the higher pressure in each
source/guest series.}\label{tab:S13}\\
\toprule
Context & Guest & Low $p$ & High $p$ & Neg.$\rightarrow$pos. &
Pos.$\rightarrow$neg. & Low/high $\rho$ \\
 & & (bar) & (bar) & fraction & fraction & \\
\midrule
\endfirsthead
\toprule
Context & Guest & Low $p$ & High $p$ & Neg.$\rightarrow$pos. &
Pos.$\rightarrow$neg. & Low/high $\rho$ \\
\midrule
\endhead
Landfill gas & CH$_4$ & 0.01 & 4.4 & 0.266 & 0.098 & 0.349 \\
Landfill gas & CO$_2$ & 0.26 & 3.2 & 0.210 & 0.103 & 0.443 \\
Methane storage & CH$_4$ & 5.8 & 65 & 0.268 & 0.058 & 0.289 \\
Natural-gas purification & CH$_4$ & 0.9 & 9 & 0.268 & 0.077 & 0.329 \\
Natural-gas purification & CO$_2$ & 0.1 & 1 & 0.191 & 0.099 & 0.526 \\
Post-combustion & CO$_2$ & 0.015 & 0.15 & 0.116 & 0.122 & 0.715 \\
Post-combustion & N$_2$ & 0.075 & 0.75 & 0.180 & 0.115 & 0.623 \\
Pre-combustion & CO$_2$ & 0.4 & 16 & 0.435 & 0.098 & -0.253 \\
Pre-combustion & H$_2$ & 0.6 & 24 & 0.150 & 0.179 & 0.585 \\
\bottomrule
\end{longtable}
\endgroup

Across individual pairs, the median fraction of positive endpoints was 0.444 at
lower pressure and 0.556 at higher pressure. Negative-to-positive transitions
substantially exceed the reverse transition for methane storage, landfill CH$_4$,
natural-gas-purification CH$_4$, and especially pre-combustion CO$_2$. These
fractions motivate the archetype analysis in \cref{sec:archetypes} but do not
identify a continuous crossover pressure or kinetic transition.

\section{Response archetypes}
\label{sec:archetypes}

Overall archetype counts and their geometry/topology composition are reported in
\cref{tab:S14,tab:S15,tab:S16}.


A pair was classified as majority positive within a pressure regime when strictly
more than half of its observed endpoint effects were positive. This generated four
deterministic archetypes.

\begin{table}[H]
\centering
\caption{Low/high-pressure response archetype census.}
\label{tab:S14}
\begin{tabular}{lrr}
\toprule
Archetype & Pairs & Fraction \\
\midrule
Positive both & 3,618 & 0.3268 \\
Pressure activated & 3,100 & 0.2800 \\
Pressure suppressed & 687 & 0.0620 \\
Negative both & 3,667 & 0.3312 \\
\bottomrule
\end{tabular}
\end{table}

\begingroup
\small
\setlength{\tabcolsep}{3pt}
\begin{longtable}{lrrrrr}
\caption{Response-archetype composition by frozen largest-contrast geometry descriptor.}
\label{tab:S15}\\
\toprule
Class & Pairs & Positive both & Activated & Suppressed & Negative both \\
\midrule
\endfirsthead
\toprule
Class & Pairs & Positive both & Activated & Suppressed & Negative both \\
\midrule
\endhead
AVA & 3,933 & 0.2929 & 0.4038 & 0.0290 & 0.2743 \\
$D_f$ & 2,211 & 0.2501 & 0.1072 & 0.1212 & 0.5215 \\
ASA & 2,036 & 0.5737 & 0.3492 & 0.0250 & 0.0521 \\
AVAf & 1,580 & 0.2297 & 0.2513 & 0.0778 & 0.4411 \\
$D_i$ & 556 & 0.2428 & 0.1097 & 0.0917 & 0.5558 \\
$D_{\mathrm{if}}$ & 540 & 0.2426 & 0.1259 & 0.1241 & 0.5074 \\
POAVA & 156 & 0.5577 & 0.2115 & 0.0577 & 0.1731 \\
POAVAf & 60 & 0.4833 & 0.0833 & 0.0667 & 0.3667 \\
\bottomrule
\end{longtable}
\endgroup

\begin{table}[H]
\centering
\caption{Response archetypes in the four most represented topology classes.}
\label{tab:S16}
\begin{tabular}{lrrrrr}
\toprule
Topology & Pairs & Positive both & Pressure activated & Pressure suppressed &
Negative both \\
\midrule
pcu & 4,662 & 0.2602 & 0.2810 & 0.0547 & 0.4041 \\
fsc & 4,219 & 0.3638 & 0.3354 & 0.0581 & 0.2427 \\
nbo & 1,055 & 0.3791 & 0.1242 & 0.0976 & 0.3991 \\
sra & 653 & 0.4686 & 0.1853 & 0.0536 & 0.2925 \\
\bottomrule
\end{tabular}
\end{table}


The four regions in main-text Figure 4A are direct visualizations of the
deterministic majority-sign definition in \cref{tab:S14}. \Cref{tab:S15} shows that
ASA/AVA-labelled pairs are enriched in positive-both or pressure-activated regimes,
whereas diameter-labelled classes are enriched in negative-both behavior.
\Cref{tab:S16} shows that topology changes the prevalence of activation without
identifying one topology as a universal mechanism.

\section{Dominant-context robustness}
\label{sec:robustness}

Class/topology removal and non-ASA/AVA contrast-tier checks are summarized in
\cref{tab:S17,tab:S18}.


Robustness analyses were performed only after completion of the confirmatory
analysis. They preserved the frozen pair identities, pair orientation, 10\%
primary geometry threshold, confirmatory endpoint, and median family-level
estimand. The original confirmatory confidence interval was not replaced.

\begin{table}[H]
\centering
\caption{Dominant-context stress tests for the primary CH$_4$ endpoint.}
\label{tab:S17}
\begin{tabular}{lrrr}
\toprule
Population & Pairs retained & Median $\Delta q$ & Positive fraction \\
 & & (\si{\milli\mole\per\gram}) & \\
\midrule
Full cohort & 11,072 & 0.5836 & 0.7050 \\
Exclude pcu & 6,410 & 0.4778 & 0.7092 \\
Exclude fsc & 6,853 & 0.5841 & 0.6822 \\
Exclude pcu + fsc & 2,191 & 0.2953 & 0.6458 \\
Exclude ASA-labelled pairs & 9,036 & 0.3642 & 0.6546 \\
Exclude AVA-labelled pairs & 7,139 & 0.3120 & 0.6296 \\
Exclude ASA + AVA-labelled pairs & 5,103 & 0.0185 & 0.5103 \\
\bottomrule
\end{tabular}
\end{table}

For the aggressive topology stress test excluding both pcu and fsc, a descriptive
10,000-replicate family bootstrap gave a 95\% interval of
\SIrange{0.2491}{0.3505}{\milli\mole\per\gram}. For the aggressive descriptor-class stress
test excluding both ASA and AVA labels, the corresponding interval was
\SIrange{-0.0080}{0.0415}{\milli\mole\per\gram}. These intervals are secondary
robustness summaries and do not replace the frozen confirmatory interval.

\begin{table}[H]
\centering
\caption{Stronger-contrast robustness after simultaneous removal of ASA- and
AVA-labelled pairs.}
\label{tab:S18}
\begin{tabular}{lrrr}
\toprule
Geometry tier & Pairs retained & Median $\Delta q$ & Positive fraction \\
 & & (\si{\milli\mole\per\gram}) & \\
\midrule
$\geq10\%$ & 5,103 & 0.0185 & 0.5103 \\
$\geq20\%$ & 2,384 & 0.1678 & 0.5600 \\
$\geq30\%$ & 1,353 & 0.7182 & 0.6526 \\
\bottomrule
\end{tabular}
\end{table}

The topology stress tests show that the pooled positive high-pressure response is
not created solely by pcu or fsc. The largest-contrast-class stress tests are more restrictive:
the near-zero response after simultaneous removal of ASA and AVA demonstrates that
the primary 10\% effect is concentrated in accessible-area and accessible-volume
changes. The recovery of a positive response at the 20 and 30\% tiers shows,
however, that sufficiently large changes in other geometry descriptors can still
produce favorable high-pressure methane responses.

These two stress tests answer different questions. Removing pcu/fsc tests whether
the primary sign depends on the dominant recorded topologies; removing ASA/AVA
tests how much of the 10\% pooled signal is compositionally concentrated in
accessible-space classes. Neither operation changes the frozen pair catalogue or
replaces the primary estimate.


\section{Database-provenance robustness}
\label{sec:provenance}

Broad provenance, leave-source-out, and ASA/AVA provenance results are reported in
\cref{tab:S19,tab:S20,tab:S21}.


Source provenance was recovered from the ARC--MOF \texttt{DBx-} framework prefix.
DB12 and DB14 were classified as experimental-source collections; the remaining
ARC--MOF source labels present in the frozen mapping were classified as
hypothetical/computationally generated sources.

\begin{table}[H]
\centering
\caption{Broad database-provenance summary for the frozen primary CH$_4$ cohort.}
\label{tab:S19}
\begin{tabular}{lrrr}
\toprule
Population & Pairs & Median $\Delta q$ & Positive fraction \\
 & & (\si{\milli\mole\per\gram}) & \\
\midrule
Full cohort & 11,072 & 0.5836 & 0.7050 \\
Same source database & 10,844 & 0.5639 & 0.7009 \\
Cross-source & 228 & 4.7091 & 0.8991 \\
Hypothetical--hypothetical & 10,981 & 0.5854 & 0.7044 \\
Experimental--experimental & 82 & 0.3552 & 0.7805 \\
Mixed hypothetical--experimental & 9 & 6.6160 & 0.7778 \\
\bottomrule
\end{tabular}
\end{table}

The cross-source and mixed-provenance populations are small and are reported for
completeness rather than as evidence for a causal source effect. The experimental
subset is directionally consistent with the primary result but is too small to
establish experimental transferability.

\begin{table}[H]
\centering
\caption{Leave-one-major-source-out primary robustness. A source was considered
major when at least 100 frozen pairs contained one or more frameworks from that
source.}
\label{tab:S20}
\begin{tabular}{lrrrr}
\toprule
Excluded source & Pairs touching source & Pairs retained &
Median $\Delta q$ & Positive fraction \\
 & & & (\si{\milli\mole\per\gram}) & \\
\midrule
DB0 & 8,261 & 2,811 & 8.7206 & 0.8399 \\
DB5 & 2,002 & 9,070 & 0.3083 & 0.6431 \\
DB6 & 482 & 10,590 & 0.6458 & 0.7212 \\
DB13 & 235 & 10,837 & 0.5862 & 0.7049 \\
DB1 & 103 & 10,969 & 0.5879 & 0.7050 \\
\bottomrule
\end{tabular}
\end{table}

The large change in magnitude after DB0 or DB5 removal demonstrates substantial
between-source heterogeneity. It does not imply that source identity causes the
response, because database source is correlated with structure-generation protocol
and sampled geometry. The relevant robustness result is narrower: every
leave-one-major-source-out population retains a positive median primary response.

\begin{table}[H]
\centering
\caption{Database-provenance robustness of the ASA/AVA-labelled population.}
\label{tab:S21}
\begin{tabular}{lrrr}
\toprule
Population & Pairs & Median $\Delta q$ & Positive fraction \\
 & & (\si{\milli\mole\per\gram}) & \\
\midrule
All ASA/AVA pairs & 5,969 & 1.7173 & 0.8715 \\
Same source database & 5,795 & 1.6840 & 0.8699 \\
Hypothetical--hypothetical & 5,899 & 1.7441 & 0.8727 \\
Experimental--experimental & 63 & 0.3808 & 0.7778 \\
Exclude DB0 & 1,641 & 10.0976 & 0.9677 \\
Exclude DB5 & 4,513 & 0.9842 & 0.8329 \\
\bottomrule
\end{tabular}
\end{table}

Thus the accessible-area/accessible-volume signal remains strongly positive after
removal of either dominant source contribution.

The leave-source-out analysis therefore addresses whether one major source is
required for the pooled sign; it does not explain why same-source medians differ.
That second question is treated separately in the post-freeze sampled-space
comparison in \cref{sec:source-heterogeneity,tab:S32,figS:source}.

\section{Claim--evidence matrix}
\label{sec:claim-evidence}

The manuscript claim--evidence map is given in \cref{tab:S22}.


\begingroup
\small
\setlength{\tabcolsep}{4pt}
\begin{longtable}{p{0.24\textwidth}p{0.29\textwidth}p{0.34\textwidth}}
\caption{Final claim--evidence matrix and explicit boundary conditions.}
\label{tab:S22}\\
\toprule
Claim & Direct evidence & Boundary / counterexample \\
\midrule
\endfirsthead
\toprule
Claim & Direct evidence & Boundary / counterexample \\
\midrule
\endhead
High-pressure CH$_4$ response is positive in the frozen matched cohort &
11,072 pairs; median $\Delta q=0.5836$ \si{\milli\mole\per\gram};
95\% bootstrap CI 0.5546--0.6149; 70.5\% positive &
Association within recorded chemistry/topology matches; not a causal deformation
experiment. \\
\addlinespace
Response strengthens with geometry contrast &
20\% tier: median 1.6980; 30\% tier: median 3.1645
\si{\milli\mole\per\gram} &
Contrast tiers are selected from the frozen pair catalogue and do not establish a
continuous dose--response mechanism. \\
\addlinespace
ASA/AVA are the dominant favorable geometry classes &
ASA/AVA removal leaves median 0.0185
\si{\milli\mole\per\gram}; ASA- and AVA-labelled pairs are strongly positive &
Other geometry classes are not uniformly inactive; the non-ASA/AVA 30\% tier has
median 0.7182 \si{\milli\mole\per\gram}. \\
\addlinespace
Geometry response depends on pressure regime &
28.0\% pressure activated versus 6.2\% pressure suppressed; median positive
endpoint fraction rises from 44.4\% to 55.6\% &
Based on discrete available low/high endpoints, not a continuously resolved
transition pressure. \\
\addlinespace
Capacity gains do not imply separation gains &
Methane-storage and pre-combustion working capacity improve on median, while
selectivity decreases in all four separation contexts &
Process metrics are equilibrium screening metrics, not complete cyclic process
simulations. \\
\addlinespace
Primary sign is not created by dominant topologies &
Remove pcu+fsc: 2,191 pairs; median 0.2953
\si{\milli\mole\per\gram}; 64.6\% positive &
Topology still modifies effect magnitude and archetype prevalence. \\
\addlinespace
Primary and ASA/AVA signs are not created by one ARC--MOF source &
Every leave-one-major-source-out primary population remains positive; ASA/AVA
remains positive after DB0 or DB5 removal &
99.2\% of frozen pairs are hypothetical--hypothetical; experimental transfer is
weakly sampled. \\
\addlinespace
Accessible-space/diameter separation survives cleaner attribution &
At top-1/top-2 $\geq1.50$, ASA remains 90.7\% positive with median 1.0699
\si{\milli\mole\per\gram}, while $D_f$ remains slightly negative &
The strict pooled cohort changes composition and becomes strongly $D_f$-weighted;
the label remains an attribution class, not an isolated causal variable. \\
\addlinespace
Positive high-pressure response survives a change of uptake basis but is density
coupled &
v/v median $+3.7972$; density-change correlations are $\rho=-0.7720$ for
\si{\milli\mole\per\gram} and $-0.3254$ for v/v &
Crystallographic v/v is not packed-bed/device performance, and density is a coupled
structural variable rather than an independent causal knob. \\
\addlinespace
Structural exemplars support a regime interpretation without mechanistic
localization &
Deterministic Figure 5 cases span favorable accessible-space, unfavorable or
pressure-suppressed diameter, and experimental-source structural contexts &
Renderings do not identify adsorption sites, measured bottlenecks, or experimental
adsorption validation. \\
\bottomrule
\end{longtable}
\endgroup


The complete frozen secondary adsorption and process outputs, together with the
source/code map, are provided in \cref{tab:S23,tab:S24,tab:S25}.

% Generated from the complete audited Step 3.5 secondary-summary CSVs.
% Project17_SI_closure_tables.tex is the frozen generated closure-table include.
% Values in that generated file are not edited here; only page orientation and
% typography are adjusted so the complete tables remain readable.
\clearpage
\begin{landscape}
\begingroup
\scriptsize
\setlength{\tabcolsep}{2.5pt}
% Auto-generated by Step 3 production/14_finalize_si_reporting.py.
% This file reports already-audited results only. Do not edit values manually.

\section{Complete secondary-result reporting}

Tables~S23--S24 expose the complete frozen Step 3.5 summary outputs. The compact
main-text and earlier SI tables emphasize the principal \si{\milli\mole\per\gram}
and process findings, whereas these tables report every audited secondary metric
without outcome-based selection.

\begin{longtable}{llrrlrrrr}
\caption{Complete frozen secondary adsorption summary. All rows from the audited \texttt{secondary\_adsorption\_summary.csv} are reported; no endpoint is selected or omitted on the basis of effect magnitude.}\label{tab:S23}\\
\toprule
Context & Guest & $T$ (K) & $p$ (bar) & Metric & Pairs & Median $\Delta$ & Positive fraction & Median sym. rel. \\
\midrule
\endfirsthead
\toprule
Context & Guest & $T$ (K) & $p$ (bar) & Metric & Pairs & Median $\Delta$ & Positive fraction & Median sym. rel. \\
\midrule
\endhead
landfill-gas & CH$_4$ & 338 & 0.01 & S(g1) & 11,011 & 0.30757 & 0.6857 & 0.24927 \\
landfill-gas & CH$_4$ & 338 & 0.01 & HOA (\si{\kilo\calorie\per\mole}) & 11,016 & -0.16007 & 0.2584 & -0.04216 \\
landfill-gas & CH$_4$ & 338 & 0.01 & \si{\milli\mole\per\gram} & 11,072 & -0.00039 & 0.4417 & -0.02747 \\
landfill-gas & CH$_4$ & 338 & 0.01 & v/v & 11,072 & -0.01244 & 0.3786 & -0.06438 \\
landfill-gas & CH$_4$ & 338 & 0.01 & wt\% & 11,072 & -0.00062 & 0.4417 & -0.02749 \\
landfill-gas & CH$_4$ & 338 & 4.4 & S(g1) & 11,045 & 0.02605 & 0.685 & 0.10246 \\
landfill-gas & CH$_4$ & 338 & 4.4 & HOA (\si{\kilo\calorie\per\mole}) & 11,046 & -0.22702 & 0.3431 & -0.06199 \\
landfill-gas & CH$_4$ & 338 & 4.4 & \si{\milli\mole\per\gram} & 11,072 & 0.06282 & 0.6096 & 0.0438 \\
landfill-gas & CH$_4$ & 338 & 4.4 & v/v & 11,072 & 0.07416 & 0.5104 & 0.00397 \\
landfill-gas & CH$_4$ & 338 & 4.4 & wt\% & 11,072 & 0.10051 & 0.6096 & 0.0438 \\
landfill-gas & CO$_2$ & 338 & 0.26 & S(g1) & 11,072 & -0.12626 & 0.3181 & -0.24627 \\
landfill-gas & CO$_2$ & 338 & 0.26 & HOA (\si{\kilo\calorie\per\mole}) & 11,042 & -0.31946 & 0.3026 & -0.05875 \\
landfill-gas & CO$_2$ & 338 & 0.26 & \si{\milli\mole\per\gram} & 11,072 & -0.06156 & 0.3306 & -0.12812 \\
landfill-gas & CO$_2$ & 338 & 0.26 & v/v & 11,072 & -1.26706 & 0.3066 & -0.16112 \\
landfill-gas & CO$_2$ & 338 & 0.26 & wt\% & 11,072 & -0.27091 & 0.3306 & -0.12811 \\
landfill-gas & CO$_2$ & 338 & 3.2 & S(g1) & 11,072 & -0.36056 & 0.3167 & -0.10178 \\
landfill-gas & CO$_2$ & 338 & 3.2 & HOA (\si{\kilo\calorie\per\mole}) & 11,052 & -0.36072 & 0.3477 & -0.06701 \\
landfill-gas & CO$_2$ & 338 & 3.2 & \si{\milli\mole\per\gram} & 11,072 & -0.09614 & 0.4378 & -0.02543 \\
landfill-gas & CO$_2$ & 338 & 3.2 & v/v & 11,072 & -3.57258 & 0.3624 & -0.06213 \\
landfill-gas & CO$_2$ & 338 & 3.2 & wt\% & 11,072 & -0.42312 & 0.4378 & -0.02543 \\
methane-storage & CH$_4$ & 298 & 5.8 & HOA (\si{\kilo\calorie\per\mole}) & 11,048 & -0.1556 & 0.272 & -0.04061 \\
methane-storage & CH$_4$ & 298 & 5.8 & \si{\milli\mole\per\gram} & 11,072 & -0.00465 & 0.4947 & -0.00114 \\
methane-storage & CH$_4$ & 298 & 5.8 & v/v & 11,072 & -1.87447 & 0.4022 & -0.03412 \\
methane-storage & CH$_4$ & 298 & 5.8 & wt\% & 11,072 & -0.00744 & 0.4947 & -0.00114 \\
methane-storage & CH$_4$ & 298 & 65 & HOA (\si{\kilo\calorie\per\mole}) & 11,050 & -0.1233 & 0.3321 & -0.03106 \\
methane-storage & CH$_4$ & 298 & 65 & \si{\milli\mole\per\gram} & 11,072 & 0.58355 & 0.705 & 0.04896 \\
methane-storage & CH$_4$ & 298 & 65 & v/v & 11,072 & 3.79719 & 0.6088 & 0.01908 \\
methane-storage & CH$_4$ & 298 & 65 & wt\% & 11,072 & 0.93369 & 0.705 & 0.04896 \\
natural-gas-purification & CH$_4$ & 298 & 0.9 & S(g1) & 11,039 & 0.03735 & 0.686 & 0.1796 \\
natural-gas-purification & CH$_4$ & 298 & 0.9 & HOA (\si{\kilo\calorie\per\mole}) & 11,043 & -0.18314 & 0.3027 & -0.05022 \\
natural-gas-purification & CH$_4$ & 298 & 0.9 & \si{\milli\mole\per\gram} & 11,072 & -0.02882 & 0.4226 & -0.03658 \\
natural-gas-purification & CH$_4$ & 298 & 0.9 & v/v & 11,072 & -0.73442 & 0.3713 & -0.06604 \\
natural-gas-purification & CH$_4$ & 298 & 0.9 & wt\% & 11,072 & -0.04611 & 0.4226 & -0.03658 \\
natural-gas-purification & CH$_4$ & 298 & 9 & S(g1) & 11,044 & 0.01961 & 0.6746 & 0.12074 \\
natural-gas-purification & CH$_4$ & 298 & 9 & HOA (\si{\kilo\calorie\per\mole}) & 11,047 & -0.23255 & 0.3606 & -0.06521 \\
natural-gas-purification & CH$_4$ & 298 & 9 & \si{\milli\mole\per\gram} & 11,072 & 0.19333 & 0.6143 & 0.04423 \\
natural-gas-purification & CH$_4$ & 298 & 9 & v/v & 11,072 & 0.49709 & 0.5216 & 0.00765 \\
natural-gas-purification & CH$_4$ & 298 & 9 & wt\% & 11,072 & 0.30932 & 0.6143 & 0.04423 \\
natural-gas-purification & CO$_2$ & 298 & 0.1 & S(g1) & 11,072 & -0.68615 & 0.316 & -0.17819 \\
natural-gas-purification & CO$_2$ & 298 & 0.1 & HOA (\si{\kilo\calorie\per\mole}) & 11,045 & -0.32475 & 0.3287 & -0.05851 \\
natural-gas-purification & CO$_2$ & 298 & 0.1 & \si{\milli\mole\per\gram} & 11,072 & -0.07351 & 0.3268 & -0.14993 \\
natural-gas-purification & CO$_2$ & 298 & 0.1 & v/v & 11,072 & -1.52278 & 0.3049 & -0.18142 \\
natural-gas-purification & CO$_2$ & 298 & 0.1 & wt\% & 11,072 & -0.32354 & 0.3268 & -0.14993 \\
natural-gas-purification & CO$_2$ & 298 & 1 & S(g1) & 11,072 & -0.63521 & 0.3271 & -0.12024 \\
natural-gas-purification & CO$_2$ & 298 & 1 & HOA (\si{\kilo\calorie\per\mole}) & 11,047 & -0.40898 & 0.3646 & -0.07507 \\
natural-gas-purification & CO$_2$ & 298 & 1 & \si{\milli\mole\per\gram} & 11,072 & -0.11967 & 0.4184 & -0.04507 \\
natural-gas-purification & CO$_2$ & 298 & 1 & v/v & 11,072 & -3.57423 & 0.3568 & -0.08329 \\
natural-gas-purification & CO$_2$ & 298 & 1 & wt\% & 11,072 & -0.52665 & 0.4184 & -0.04507 \\
post-combustion-vsa & CO$_2$ & 298 & 0.015 & S(g1) & 11,072 & -0.86314 & 0.319 & -0.29117 \\
post-combustion-vsa & CO$_2$ & 298 & 0.015 & HOA (\si{\kilo\calorie\per\mole}) & 11,043 & -0.32293 & 0.3164 & -0.05692 \\
post-combustion-vsa & CO$_2$ & 298 & 0.015 & \si{\milli\mole\per\gram} & 11,072 & -0.01793 & 0.3293 & -0.17225 \\
post-combustion-vsa & CO$_2$ & 298 & 0.015 & v/v & 11,072 & -0.36175 & 0.3046 & -0.20939 \\
post-combustion-vsa & CO$_2$ & 298 & 0.015 & wt\% & 11,072 & -0.0789 & 0.3293 & -0.17226 \\
post-combustion-vsa & CO$_2$ & 298 & 0.15 & S(g1) & 11,072 & -4.32004 & 0.3142 & -0.23387 \\
post-combustion-vsa & CO$_2$ & 298 & 0.15 & HOA (\si{\kilo\calorie\per\mole}) & 11,049 & -0.33993 & 0.3074 & -0.06081 \\
post-combustion-vsa & CO$_2$ & 298 & 0.15 & \si{\milli\mole\per\gram} & 11,072 & -0.11955 & 0.3235 & -0.14864 \\
post-combustion-vsa & CO$_2$ & 298 & 0.15 & v/v & 11,072 & -2.40602 & 0.3036 & -0.17904 \\
post-combustion-vsa & CO$_2$ & 298 & 0.15 & wt\% & 11,072 & -0.52614 & 0.3235 & -0.14864 \\
post-combustion-vsa & N$_2$ & 298 & 0.075 & S(g1) & 11,043 & 0.05072 & 0.6828 & 0.2941 \\
post-combustion-vsa & N$_2$ & 298 & 0.075 & HOA (\si{\kilo\calorie\per\mole}) & 11,070 & -0.11165 & 0.2558 & -0.04809 \\
post-combustion-vsa & N$_2$ & 298 & 0.075 & \si{\milli\mole\per\gram} & 11,072 & 0.00038 & 0.5424 & 0.0125 \\
post-combustion-vsa & N$_2$ & 298 & 0.075 & v/v & 11,072 & -0.02662 & 0.3167 & -0.05566 \\
post-combustion-vsa & N$_2$ & 298 & 0.075 & wt\% & 11,072 & 0.00107 & 0.5424 & 0.01252 \\
post-combustion-vsa & N$_2$ & 298 & 0.75 & S(g1) & 11,049 & 0.00616 & 0.6872 & 0.2347 \\
post-combustion-vsa & N$_2$ & 298 & 0.75 & HOA (\si{\kilo\calorie\per\mole}) & 11,072 & -0.12743 & 0.2671 & -0.05391 \\
post-combustion-vsa & N$_2$ & 298 & 0.75 & \si{\milli\mole\per\gram} & 11,072 & 0.00431 & 0.6075 & 0.02933 \\
post-combustion-vsa & N$_2$ & 298 & 0.75 & v/v & 11,072 & -0.06395 & 0.3679 & -0.02876 \\
post-combustion-vsa & N$_2$ & 298 & 0.75 & wt\% & 11,072 & 0.01207 & 0.6075 & 0.02933 \\
pre-combustion-40-40 & CO$_2$ & 313 & 0.4 & S(g1) & 11,072 & -7.82344 & 0.2864 & -0.26498 \\
pre-combustion-40-40 & CO$_2$ & 313 & 0.4 & HOA (\si{\kilo\calorie\per\mole}) & 11,050 & -0.33278 & 0.2982 & -0.06108 \\
pre-combustion-40-40 & CO$_2$ & 313 & 0.4 & \si{\milli\mole\per\gram} & 11,072 & -0.15609 & 0.332 & -0.12163 \\
pre-combustion-40-40 & CO$_2$ & 313 & 0.4 & v/v & 11,072 & -3.14037 & 0.3096 & -0.15229 \\
pre-combustion-40-40 & CO$_2$ & 313 & 0.4 & wt\% & 11,072 & -0.68693 & 0.332 & -0.12163 \\
pre-combustion-40-40 & CO$_2$ & 313 & 16 & S(g1) & 11,071 & -4.5636 & 0.3821 & -0.06789 \\
pre-combustion-40-40 & CO$_2$ & 313 & 16 & HOA (\si{\kilo\calorie\per\mole}) & 11,058 & -0.38757 & 0.3621 & -0.0631 \\
pre-combustion-40-40 & CO$_2$ & 313 & 16 & \si{\milli\mole\per\gram} & 11,071 & 0.54859 & 0.669 & 0.04356 \\
pre-combustion-40-40 & CO$_2$ & 313 & 16 & v/v & 11,071 & 4.61534 & 0.5934 & 0.02079 \\
pre-combustion-40-40 & CO$_2$ & 313 & 16 & wt\% & 11,071 & 2.41431 & 0.669 & 0.04356 \\
pre-combustion-40-40 & H$_2$ & 313 & 0.6 & S(g1) & 11,050 & 0.00442 & 0.715 & 0.26631 \\
pre-combustion-40-40 & H$_2$ & 313 & 0.6 & HOA (\si{\kilo\calorie\per\mole}) & 11,071 & -0.0583 & 0.2239 & -0.03917 \\
pre-combustion-40-40 & H$_2$ & 313 & 0.6 & \si{\milli\mole\per\gram} & 11,072 & 0.00262 & 0.7338 & 0.05239 \\
pre-combustion-40-40 & H$_2$ & 313 & 0.6 & v/v & 11,072 & 0.02718 & 0.6787 & 0.03133 \\
pre-combustion-40-40 & H$_2$ & 313 & 0.6 & wt\% & 11,072 & 0.00052 & 0.7337 & 0.05239 \\
pre-combustion-40-40 & H$_2$ & 313 & 24 & S(g1) & 11,058 & 0.00067 & 0.6186 & 0.06814 \\
pre-combustion-40-40 & H$_2$ & 313 & 24 & HOA (\si{\kilo\calorie\per\mole}) & 11,071 & -0.10048 & 0.3798 & -0.05515 \\
pre-combustion-40-40 & H$_2$ & 313 & 24 & \si{\milli\mole\per\gram} & 11,071 & 0.02185 & 0.7054 & 0.09629 \\
pre-combustion-40-40 & H$_2$ & 313 & 24 & v/v & 11,071 & 0.30484 & 0.6904 & 0.07893 \\
pre-combustion-40-40 & H$_2$ & 313 & 24 & wt\% & 11,071 & 0.00437 & 0.7054 & 0.0963 \\
\bottomrule
\end{longtable}

\begin{longtable}{llrrrr}
\caption{Complete frozen secondary process summary. All rows from the audited \texttt{secondary\_process\_summary.csv} are reported.}\label{tab:S24}\\
\toprule
Process & Metric & Pairs & Median $\Delta$ & Positive fraction & Median sym. rel. \\
\midrule
\endfirsthead
\toprule
Process & Metric & Pairs & Median $\Delta$ & Positive fraction & Median sym. rel. \\
\midrule
\endhead
landfill-gas-vpsa & AFM & 11,008 & -0.6251 & 0.4947 & -0.00915 \\
landfill-gas-vpsa & Uptake (\si{\milli\mole\per\gram}) & 11,008 & -0.10031 & 0.4346 & -0.02658 \\
landfill-gas-vpsa & Working capacity (\si{\milli\mole\per\gram}) & 11,008 & -0.02488 & 0.4836 & -0.0086 \\
landfill-gas-vpsa & Purity & 11,008 & -0.01656 & 0.3527 & -0.02367 \\
landfill-gas-vpsa & Selectivity & 11,008 & -0.363 & 0.3147 & -0.10228 \\
landfill-gas-vpsa & SSP & 11,008 & -1.36634 & 0.3328 & -0.17922 \\
landfill-gas-vpsa & Uptake (v/v) & 11,008 & -3.67014 & 0.3588 & -0.06379 \\
landfill-gas-vpsa & Working capacity (v/v) & 11,008 & -1.84331 & 0.4142 & -0.04199 \\
landfill-gas-vpsa & Uptake (wt\%) & 11,008 & -0.44146 & 0.4346 & -0.02658 \\
landfill-gas-vpsa & Working capacity (wt\%) & 11,008 & -0.1095 & 0.4836 & -0.0086 \\
methane-storage-psa & AFM & 0 & -- & -- & -- \\
methane-storage-psa & Uptake (\si{\milli\mole\per\gram}) & 11,008 & 0.57487 & 0.7033 & 0.04827 \\
methane-storage-psa & Working capacity (\si{\milli\mole\per\gram}) & 11,008 & 0.46684 & 0.7231 & 0.07219 \\
methane-storage-psa & Purity & 0 & -- & -- & -- \\
methane-storage-psa & Selectivity & 0 & -- & -- & -- \\
methane-storage-psa & SSP & 0 & -- & -- & -- \\
methane-storage-psa & Uptake (v/v) & 11,008 & 3.67548 & 0.6066 & 0.01862 \\
methane-storage-psa & Working capacity (v/v) & 11,008 & 5.69722 & 0.6893 & 0.0512 \\
methane-storage-psa & Uptake (wt\%) & 11,008 & 0.9198 & 0.7033 & 0.04827 \\
methane-storage-psa & Working capacity (wt\%) & 11,008 & 0.74695 & 0.7231 & 0.07219 \\
natural-gas-purification & AFM & 11,008 & -1.42043 & 0.4366 & -0.10071 \\
natural-gas-purification & Uptake (\si{\milli\mole\per\gram}) & 11,008 & -0.12502 & 0.4152 & -0.04664 \\
natural-gas-purification & Working capacity (\si{\milli\mole\per\gram}) & 11,008 & -0.02018 & 0.4849 & -0.00961 \\
natural-gas-purification & Purity & 11,008 & -0.02733 & 0.342 & -0.07405 \\
natural-gas-purification & Selectivity & 11,008 & -0.63771 & 0.3256 & -0.12048 \\
natural-gas-purification & SSP & 11,008 & -0.70277 & 0.3352 & -0.24103 \\
natural-gas-purification & Uptake (v/v) & 11,008 & -3.68349 & 0.3533 & -0.08541 \\
natural-gas-purification & Working capacity (v/v) & 11,008 & -1.28715 & 0.4263 & -0.04683 \\
natural-gas-purification & Uptake (wt\%) & 11,008 & -0.5502 & 0.4152 & -0.04664 \\
natural-gas-purification & Working capacity (wt\%) & 11,008 & -0.08881 & 0.4849 & -0.00961 \\
post-combustion-vsa & AFM & 11,008 & -7.84541 & 0.3705 & -0.38243 \\
post-combustion-vsa & Uptake (\si{\milli\mole\per\gram}) & 11,008 & -0.12275 & 0.3197 & -0.15156 \\
post-combustion-vsa & Working capacity (\si{\milli\mole\per\gram}) & 11,008 & -0.10279 & 0.3308 & -0.14981 \\
post-combustion-vsa & Purity & 11,008 & -0.01526 & 0.3185 & -0.01785 \\
post-combustion-vsa & Selectivity & 11,008 & -4.45592 & 0.3106 & -0.2367 \\
post-combustion-vsa & SSP & 11,008 & -35.10245 & 0.315 & -0.41811 \\
post-combustion-vsa & Uptake (v/v) & 11,008 & -2.46513 & 0.3 & -0.18114 \\
post-combustion-vsa & Working capacity (v/v) & 11,008 & -2.03925 & 0.31 & -0.18179 \\
post-combustion-vsa & Uptake (wt\%) & 11,008 & -0.54023 & 0.3197 & -0.15156 \\
post-combustion-vsa & Working capacity (wt\%) & 11,008 & -0.45239 & 0.3308 & -0.14981 \\
pre-combustion-40-40 & AFM & 11,008 & 417.47241 & 0.6578 & 0.26756 \\
pre-combustion-40-40 & Uptake (\si{\milli\mole\per\gram}) & 11,008 & 0.53849 & 0.6672 & 0.04268 \\
pre-combustion-40-40 & Working capacity (\si{\milli\mole\per\gram}) & 11,008 & 0.66027 & 0.6921 & 0.06763 \\
pre-combustion-40-40 & Purity & 11,008 & -0.00043 & 0.4584 & -0.00043 \\
pre-combustion-40-40 & Selectivity & 11,008 & -4.70478 & 0.3786 & -0.06934 \\
pre-combustion-40-40 & SSP & 11,008 & -155.38651 & 0.4276 & -0.09183 \\
pre-combustion-40-40 & Uptake (v/v) & 11,008 & 4.46127 & 0.591 & 0.02014 \\
pre-combustion-40-40 & Working capacity (v/v) & 11,008 & 7.3928 & 0.6409 & 0.04113 \\
pre-combustion-40-40 & Uptake (wt\%) & 11,008 & 2.36987 & 0.6672 & 0.04268 \\
pre-combustion-40-40 & Working capacity (wt\%) & 11,008 & 2.90581 & 0.6921 & 0.06763 \\
\bottomrule
\end{longtable}

\section{Reproducibility map}

The analysis used ARC--MOF release v7 (Zenodo record 16802743). The local raw-data
inventory comprised the charged CIF archive together with the ARC--MOF topology,
dimensionality, geometry, RAC/cluster, adsorption, and process tables used by the
staged workflow. Each analytical stage wrote deterministic tables and a manifest;
each stage was followed by an independent reconstruction audit.

\begin{longtable}{lll}
\caption{Source-data and code map for manuscript figures and late-stage Supporting Information tables.}\label{tab:S25}\\
\toprule
Output & Frozen source data & Production script \\
\midrule
\endfirsthead
\toprule
Output & Frozen source data & Production script \\
\midrule
\endhead
Figure 1 & Stages 1.1--1.4; Steps 2.1--2.3 frozen cohort and pair tables & \texttt{03\_make\_figure1.py} \\
Figure 2 & Step 3.4 \texttt{primary\_ch4\_pair\_effects.parquet} and summaries & \texttt{05\_make\_figure2.py} \\
Figure 3 & Step 3.5 secondary adsorption/process summary tables & \texttt{07\_make\_figure3\_clean\_v2.py} \\
Figure 4 & Step 3.7 response archetype tables & \texttt{10\_make\_figure4\_v2.py} \\
Figures S1--S3 & Steps 3.5--3.6 secondary and concordance tables & \texttt{11\_make\_si\_figures.py} \\
Tables S17--S18 & Step 3.8 dominant-context robustness tables & \texttt{12\_analyze\_robustness\_transferability.py} \\
Tables S19--S21 & Step 3.10 database-provenance robustness tables & \texttt{13\_analyze\_database\_provenance.py} \\
Tables S23--S24 & Audited full Step 3.5 summary tables & \texttt{14\_finalize\_si\_reporting.py} \\
\bottomrule
\end{longtable}

\endgroup
\end{landscape}
\clearpage


\section{Post-freeze geometry-attribution stability}
\label{sec:attribution}

The pooled and descriptor-resolved dominance checks are tabulated in
\cref{tab:S26,tab:S27} and summarized graphically in \cref{figS:attribution}.


The frozen pair label is the descriptor with the largest symmetric relative geometry
difference in a selected pair. Because the eight pore descriptors can change
simultaneously, a post-freeze analysis quantified how clearly the largest difference
exceeded the second largest. No pair was reselected and no outcome informed the
dominance criterion.

\begin{table}[H]
\centering
\caption{Pooled primary CH$_4$ response after progressively requiring the largest
pairwise geometry difference to dominate the second-largest difference.}
\label{tab:S26}
\begin{tabular}{lrrrr}
\toprule
Top-1/top-2 criterion & Pairs & Fraction of frozen cohort & Median $\Delta q$ & Positive fraction \\
 & & & (\si{\milli\mole\per\gram}) & \\
\midrule
$\geq1.10$ & 4,603 & 0.4157 & 0.4135 & 0.6626 \\
$\geq1.25$ & 2,866 & 0.2589 & 0.1873 & 0.6015 \\
$\geq1.50$ & 1,715 & 0.1549 & 0.0383 & 0.5312 \\
\bottomrule
\end{tabular}
\end{table}

The pooled attenuation should not be read as a uniform loss of descriptor-specific
response because the retained class composition changes strongly as the dominance
criterion tightens. \Cref{tab:S27} therefore reports the major classes separately.

\begin{table}[H]
\centering
\caption{Descriptor-resolved primary response in increasingly clean
largest-contrast subsets. Values use the frozen pair orientation and endpoint.}
\label{tab:S27}
\begin{tabular}{llrrr}
\toprule
Class & Top-1/top-2 criterion & Pairs & Median $\Delta q$ & Positive fraction \\
 & & & (\si{\milli\mole\per\gram}) & \\
\midrule
ASA & $\geq1.00$ & 2,036 & 2.9964 & 0.9288 \\
ASA & $\geq1.10$ & 883 & 1.3745 & 0.9060 \\
ASA & $\geq1.25$ & 582 & 1.3514 & 0.9038 \\
ASA & $\geq1.50$ & 311 & 1.0699 & 0.9068 \\
\addlinespace
AVA & $\geq1.00$ & 3,933 & 1.2201 & 0.8419 \\
AVA & $\geq1.10$ & 1,236 & 2.2861 & 0.9547 \\
AVA & $\geq1.25$ & 379 & 2.2727 & 0.9815 \\
AVA & $\geq1.50$ & 23 & 1.6166 & 0.9565 \\
\addlinespace
$D_f$ & $\geq1.00$ & 2,211 & -0.1457 & 0.4120 \\
$D_f$ & $\geq1.10$ & 1,975 & -0.1193 & 0.4233 \\
$D_f$ & $\geq1.25$ & 1,666 & -0.1041 & 0.4292 \\
$D_f$ & $\geq1.50$ & 1,265 & -0.0867 & 0.4364 \\
\bottomrule
\end{tabular}
\end{table}

Taken together, \cref{tab:S26,tab:S27,figS:attribution} show that the loss of pooled
support under stricter dominance is accompanied by a major shift in class composition.
Clean ASA-dominant and AVA-dominant subsets remain strongly favorable while
clean $D_f$-dominant subsets remain slightly unfavorable. At the 1.50 pooled cut,
$D_f$ contributes 1,265 of 1,715 retained pairs. The shrinking pooled median is
therefore partly a descriptor-mixture effect. The largest-contrast label remains
noncausal because descriptors co-vary, but the accessible-space/diameter
heterogeneity survives increasingly strict dominance filtering.

\sifigure{Figure_S4_geometry_attribution_stability.pdf}{figS:attribution}
  {Geometry-attribution stability in the frozen pair catalogue}
  {Panel A, top-1/top-2 contrast-ratio distribution; panel B, Spearman correlation
   among the eight pairwise geometry contrasts; panel C, descriptor-resolved median
   primary response across increasingly strict dominance cuts; panel D, median
   primary response versus the number of descriptors changing by at least 10\%.}
  {The post-freeze attribution analysis does not reselect any pair or alter its
   orientation. The pooled response attenuates as increasingly dominant single
   contrasts are required, but descriptor-resolved behavior remains heterogeneous:
   accessible-space classes remain favorable while $D_f$ remains slightly negative.
   Descriptor correlations and the multiplicity panel emphasize that the frozen
   largest-contrast label is an attribution class within multidimensional pore
   changes, not an isolated causal perturbation.}

The central point of \cref{figS:attribution} is class-resolved stability rather than
the pooled attenuation alone. Tightening top-1/top-2 dominance changes which
descriptor classes remain in the cohort; at the strictest cut, $D_f$ accounts for
1,265 of 1,715 retained pairs. ASA and AVA nevertheless retain strongly favorable
medians in their clean-dominance subsets while $D_f$ remains slightly negative.
This is the SI evidence used to bound the largest-contrast terminology in main-text
Section 2.3 and the structural interpretation of Figure 5.

\section{Post-freeze uptake-basis and density robustness}
\label{sec:basis-density}

Uptake-basis, density-correlation, density-quintile, and descriptor-group results are
reported in \cref{tab:S28,tab:S29,tab:S30,tab:S31} and visualized in
\cref{figS:basis-density}.


The 65-bar methane analysis was repeated on the same frozen pairs using the source
v/v and wt\% fields after exact reproduction of the frozen
\si{\milli\mole\per\gram} pair effects.

\begin{table}[H]
\centering
\caption{Primary 298 K, 65 bar CH$_4$ response across uptake bases.}
\label{tab:S28}
\begin{tabular}{lrrrrr}
\toprule
Metric & Pairs & Median $\Delta$ & Mean $\Delta$ & Positive fraction & 95\% bootstrap CI \\
\midrule
\si{\milli\mole\per\gram} & 11,072 & 0.5836 & 2.4275 & 0.7050 & 0.5546--0.6149 \\
v/v & 11,072 & 3.7972 & 13.8259 & 0.6088 & 3.4385--4.1851 \\
wt\% & 11,072 & 0.9337 & 3.8839 & 0.7050 & 0.8873--0.9839 \\
\bottomrule
\end{tabular}
\end{table}

\begin{table}[H]
\centering
\caption{Association between signed framework-density change and the primary
methane response.}
\label{tab:S29}
\begin{tabular}{lrr}
\toprule
Uptake basis & Pairs & Spearman $\rho$ with signed density change \\
\midrule
\si{\milli\mole\per\gram} & 11,072 & -0.7720 \\
v/v & 11,072 & -0.3254 \\
\bottomrule
\end{tabular}
\end{table}

\begin{table}[H]
\centering
\caption{Density-quintile robustness. Q1 contains the strongest density decreases
from the geometry-low to geometry-high member; Q5 contains density increases.}
\label{tab:S30}
\begin{tabular}{lrrrrr}
\toprule
Density quintile & Pairs & Median density change & Median $\Delta q$ & Positive $q$ & Median $\Delta$(v/v) \\
\midrule
Q1 & 2,215 & -0.6667 & 10.1770 & 1.0000 & 44.2788 \\
Q2 & 2,214 & -- & 1.8568 & 0.935 & 9.4190 \\
Q3 & 2,214 & -- & 0.2224 & 0.645 & 2.0099 \\
Q4 & 2,214 & -- & 0.0314 & 0.523 & 0.4049 \\
Q5 & 2,215 & $>0$ & -0.1598 & 0.422 & 0.0541 \\
\bottomrule
\end{tabular}
\end{table}

\begin{table}[H]
\centering
\caption{Uptake-basis response by accessible-space versus other
largest-contrast classes.}
\label{tab:S31}
\begin{tabular}{llrrr}
\toprule
Population & Basis & Pairs & Median $\Delta$ & Positive fraction \\
\midrule
ASA/AVA-labelled & \si{\milli\mole\per\gram} & 5,969 & 1.7173 & 0.8715 \\
ASA/AVA-labelled & v/v & 5,969 & 11.3730 & 0.7291 \\
Other labels & \si{\milli\mole\per\gram} & 5,103 & 0.0185 & 0.5103 \\
Other labels & v/v & 5,103 & -0.7147 & 0.4682 \\
\bottomrule
\end{tabular}
\end{table}

Across \cref{tab:S28,tab:S29,tab:S30,tab:S31,figS:basis-density}, the positive
population-level v/v median demonstrates that the high-pressure sign
is not only a gravimetric normalization artifact. At the same time, the strong
negative density correlation and the descriptor-group split show that the effect
magnitude is coupled to framework density. The v/v field is used here only as an
internal uptake-basis robustness test; it does not establish packed-bed or
system-level volumetric storage performance.

\sifigure{Figure_S5_basis_density_robustness.pdf}{figS:basis-density}
  {Uptake-basis and density robustness of the 65-bar methane response}
  {Panel A, median change on the native gravimetric, volumetric, and wt\% source
   bases; panel B, pair-level gravimetric versus volumetric methane response;
   panel C, median volumetric response across ordered density-change quintiles;
   panel D, positive-pair fraction for ASA/AVA-labelled versus other pairs on
   gravimetric and volumetric bases.}
  {The population-level response remains positive after volumetric normalization,
   while pair-level consistency decreases. Density-change stratification and the
   ASA/AVA split show that response magnitude is strongly coupled to framework
   density and accessible-space class. The v/v comparison is an internal
   crystallographic-basis robustness test and is not interpreted as packed-bed or
   device-level volumetric storage performance.}

\Cref{figS:basis-density} separates two statements that should not be conflated.
First, the positive population sign survives when the same pair comparison is
expressed on the source crystallographic volumetric basis. Second, effect magnitude
remains strongly organized by density change and by whether the largest contrast is
ASA/AVA. The H3 analysis therefore supports a coupled pore-space/density regime; it
does not establish a density-independent geometry effect and does not justify
retrospectively adding density to the matching key.

\section{Post-freeze same-source heterogeneity}
\label{sec:source-heterogeneity}

The same-source effect, geometry magnitude, and descriptor composition are reported
in \cref{tab:S32} and visualized jointly in \cref{figS:source}.


\begin{table}[H]
\centering
\small
\caption{Same-source primary CH$_4$ response and geometry composition for source
populations containing at least 100 pairs.}
\label{tab:S32}
\begin{tabular}{lrrrrr}
\toprule
Source & Pairs & Median $\Delta q$ & Positive & Median max contrast & ASA/AVA \\
 & & (\si{\milli\mole\per\gram}) & fraction & & fraction \\
\midrule
DB0 & 8,053 & 0.3420 & 0.6528 & 0.2253 & 0.5181 \\
DB5 & 1,808 & 10.4372 & 0.9945 & 0.9580 & 0.7218 \\
DB6 & 476 & -0.2445 & 0.3424 & 0.1802 & 0.0798 \\
DB13 & 225 & 0.4664 & 0.7022 & 0.1722 & 0.3911 \\
\bottomrule
\end{tabular}
\end{table}

As shown jointly by \cref{tab:S32,figS:source}, the source strata occupy very
different geometry regimes. This is an applicability
diagnostic, not evidence that database provenance itself causes the adsorption
response.

\sifigure{Figure_S6_source_stratified.pdf}{figS:source}
  {Same-source heterogeneity in the frozen primary methane response}
  {Same-source populations containing at least 100 frozen pairs. Panel A, median
   298 K, 65 bar CH$_4$ response; panel B, median maximum frozen geometry contrast;
   panel C, fraction of pairs labelled by ASA or AVA as the largest-contrast class.}
  {The source-specific medians differ strongly because the source populations occupy
   different geometry regimes and descriptor compositions. The juxtaposition of
   effect magnitude, geometry contrast, and ASA/AVA composition makes this sampling
   difference explicit. Source is therefore treated as an applicability and
   provenance diagnostic rather than as a causal explanation of adsorption response.}

The juxtaposition matters quantitatively. DB5 combines a much larger median geometry
contrast (0.9580) and a 72.2\% ASA/AVA composition with a strongly positive median
response, whereas DB6 combines a smaller median contrast (0.1802) and only 8.0\%
ASA/AVA with a negative median. These covarying differences prevent a causal
database-source interpretation and explain why \cref{figS:source} is used as an
applicability diagnostic in the main text.

\section{Deterministic structural case-study set}
\label{sec:structures}

The six source cases used for main-text Figure 5 are reported in \cref{tab:S33}.


Six frozen pairs were selected for structural rendering from predeclared response
and geometry strata. Selection was deterministic, mutually exclusive across panels,
and based on representativeness within each stratum rather than maximizing the
adsorption response. These six rows are the source cases for main-text Figure 5,
which places each frozen geometry-low/geometry-high pair side by side under one
rendering convention.

\begin{landscape}
\begin{table}[H]
\centering
\scriptsize
\caption{Frozen matched-pair examples used in main-text Figure 5. The Class column preserves the frozen ordering label while noting co-dominant geometry context where relevant.}
\label{tab:S33}
\begin{tabular}{cllllrr}
\toprule
Panel & Case & Class & Topology & Archetype & Top1/Top2 & Primary $\Delta q$ \\
 & & & & & & (\si{\milli\mole\per\gram}) \\
\midrule
A & Broad accessible-space contrast & ASA label; accessibility co-dominant & pcu & Pressure activated & 1.000 & 9.0970 \\
B & AVA-led accessible space & AVA & fsc & Pressure activated & 1.183 & 2.9942 \\
C & ASA--diameter mismatch & ASA & fsc & Positive both & 1.642 & 2.6674 \\
D & Diameter counterexample & $D_f$ & pcu & Negative both & 1.221 & -0.4598 \\
E & Pressure-suppressed diameter & $D_i$/$D_{\mathrm{if}}$ co-dominant & pcu & Pressure suppressed & 1.010 & -0.1876 \\
F & Experimental-source anchor & AVA/AVAf co-dominant & dia & Positive both & 1.026 & 0.0826 \\
\bottomrule
\end{tabular}
\end{table}
\end{landscape}

The case labels in \cref{tab:S33} are synchronized with the final main-text
Figure 5 interpretation. Panel A retains ASA as the frozen ordering label but is
described as a broad accessible-space contrast because multiple accessibility
metrics are co-dominant. Panel B is AVA-led and multidimensional. Panel C is the
ASA--diameter mismatch emphasized in the main discussion: its favorable response
occurs without a parallel AVA increase and while $D_f$ decreases. Panels D and E
are deliberately retained as diameter counter-regimes, showing that a larger
characteristic diameter is not equivalent to a favorable adsorption response.
Panel F is an experimental-source crystallographic record rather than experimental
adsorption validation.

The local insets in main-text Figure 5 provide representative pore, cavity, or
framework context only; they are not adsorption-site, measured-bottleneck, or
unique-mechanism assignments. The 12 CIFs were extracted with a checksum manifest
and are used for visualization only. Their purpose is to translate the statistical
descriptor regimes into crystallographic intuition without changing the evidential
level of the analysis.

\section{ARC--MOF-specific limitations}
\label{sec:limitations}

The present workflow inherits limitations of the ARC--MOF structural and simulation
resource. Adsorption values represent equilibrium calculations for recorded
framework structures and therefore do not include synthesis success, activation,
defect populations, phase mixtures, framework flexibility, kinetic transport, or
long-term stability. Classical interaction models and assigned partial charges
introduce force-field and electrostatic uncertainty that is not propagated into
the matched-pair confidence intervals. Geometry descriptors characterize pore size
and accessibility but do not uniquely encode local adsorption-site energetics.

The available pressure grids are sparse relative to a full isotherm. Low/high
pressure classifications and response archetypes should therefore be interpreted as
endpoint-based summaries, not as measurements of a unique transition pressure.
Likewise, process metrics are screening quantities rather than full dynamic
process simulations.

Finally, the frozen cohort is strongly dominated by hypothetical structures.
Source leave-out tests protect against the narrow claim that one ARC--MOF source
database creates the observed primary sign, but they cannot establish experimental
realizability, activation success, or synthesis robustness. The crystallographic
v/v analysis is likewise an internal source-basis check rather than a packed-bed,
pellet, monolith, or tank-level claim. Figure 5F provides experimental-source
structural provenance but not experimental adsorption validation. The complete
inferential boundary is summarized in \cref{sec:boundary}; targeted measurements in
experimentally characterized matched families remain the decisive next test.

\section{Audit trail and reproducibility}
\label{sec:audit}

The local analysis used ARC--MOF release v7 (Zenodo DOI \url{https://doi.org/10.5281/zenodo.16802743}), including
the charged CIF archive; topology, dimensionality, geometry, RAC, and chemistry
cluster tables; the nine adsorption tables; and \texttt{overall\_process.csv}.
The generated complete secondary-result tables and source/code map are retained in
\cref{tab:S23,tab:S24,tab:S25}. Every stage wrote a manifest, deterministic data
tables, and an independent audit.
The audits reconstructed source hashes, target populations, chemistry and
architecture grouping keys, exclusions, geometry ranges, threshold flags, all
within-family pair candidates, deterministic pair optimality, endpoint coverage,
primary outcome linkage, the complete 10,000-replicate bootstrap, secondary effect
arithmetic, cross-condition correlation/sign matrices, and response-archetype
assignments.

The completed audit sequence verified that:
\begin{itemize}
    \item adsorption/process outcomes were excluded from chemistry, architecture,
          threshold, and pair construction;
    \item chemistry groups were frozen before architecture refinement;
    \item geometry did not enter chemistry or architecture keys;
    \item the 10\% primary geometry threshold was frozen before pair construction;
    \item exactly one deterministic pair was frozen per eligible family;
    \item the confirmatory CH$_4$ endpoint was frozen before outcome values were read;
    \item all 11,072 primary pairs had complete confirmatory outcome values;
    \item the 10,000 bootstrap replicates reproduced exactly under the frozen seed;
    \item all secondary endpoints preserved frozen pair membership and orientation;
    \item the full cross-condition matrices reproduced from pair-level effects; and
    \item all low/high-pressure archetypes reproduced exactly from the frozen
          endpoint matrix;
    \item all ARC--MOF database-source labels reconstructed exactly from frozen
          framework identifiers; and
    \item all leave-one-major-source-out primary and ASA/AVA robustness summaries
          reproduced exactly;
    \item post-freeze top-1/top-2 attribution subsets preserved the frozen pair
          catalogue and reproduced the descriptor-resolved clean-dominance summaries;
    \item uptake-basis analysis exactly reproduced the frozen gravimetric primary
          effects before calculating v/v, wt\%, and density-stratified summaries; and
    \item the 12 structural-case CIFs were extracted from the frozen catalogue with
          a checksum manifest and were used for visualization only.
\end{itemize}

\section{Interpretive boundary}
\label{sec:boundary}

The analysis supports geometry-associated differences within families matched on the
recorded chemistry representation and broad recorded architecture. It does not
prove atom-by-atom identity and is not a causal deformation experiment. The pair
direction is defined by the frozen largest-contrast geometry label, but pore
descriptors can co-vary; the post-freeze clean-dominance results
(\cref{tab:S26,tab:S27,figS:attribution}) therefore treat that label as an
attribution class rather than an isolated causal variable. Importantly, the pooled
attenuation under stricter dominance cuts is not a uniform loss of class-specific
response: ASA and AVA remain strongly positive in clean subsets, $D_f$ remains
slightly negative, and the strict pooled population becomes heavily $D_f$-weighted.

The positive high-pressure CH$_4$ median persists on the crystallographic
volumetric basis, but density change is strongly coupled to both gravimetric and
volumetric response (\cref{tab:S28,tab:S29,tab:S30,tab:S31,figS:basis-density}).
Density is therefore retained as coupled structural context rather than introduced
post hoc as an independent nuisance variable. The v/v check is internal to the
database representation and is not a packed-density or device-level storage claim.

The primary sign survives removal of the dominant topology classes and every major
ARC--MOF source, but the frozen cohort remains 99.2\%
hypothetical--hypothetical (\cref{tab:S17,tab:S19,tab:S20,tab:S21}). Same-source
effect heterogeneity is entangled with geometry magnitude and ASA/AVA composition
(\cref{tab:S32,figS:source}), so source is an applicability label rather than a
causal explanation. The structural cases in \cref{tab:S33} provide
crystallographic intuition only; their insets are not adsorption-site or measured-
bottleneck assignments, and the experimental-source case is not experimental
adsorption validation.

Finally, every gas, pressure, process, concordance, archetype, density,
source-stratified, and structural-case analysis outside the frozen 65-bar CH$_4$
primary endpoint remains secondary or post-freeze. The summary claim--evidence
matrix in \cref{tab:S22} records the corresponding boundary for each major
manuscript conclusion.


\clearpage
\bibliographystyle{unsrtnat}
\bibliography{Project17_References}

\end{document}
